\documentclass[12pt,a4paper]{article}

\usepackage[utf8]{inputenc}
\usepackage[T1]{fontenc}
\usepackage{graphicx}
\usepackage{geometry}
\geometry{margin=1in}
\usepackage{hyperref}
\usepackage{enumitem}
\usepackage{caption}
\usepackage{float}
\usepackage{xcolor}
\usepackage{fancyhdr}
\usepackage{titlesec}
\usepackage{amsmath, amssymb}
\usepackage{array}
\usepackage{parskip}
\usepackage{booktabs}
\usepackage{longtable}
\usepackage{xurl}
\usepackage{tikz}
\usetikzlibrary{arrows.meta, positioning, shapes.multipart, shapes.geometric, calc, fit, backgrounds}
\usepackage{listings}

\urlstyle{tt}
\newcommand{\code}[1]{\texttt{\url{#1}}}
% \cls: monospaced identifier that can wrap inside narrow table cells (fixes overlapping text)
\newcommand{\cls}[1]{\nolinkurl{#1}}

% ---- Team members (edit once, used everywhere) ----
\newcommand{\ma}{Ahmad Nawaz}
\newcommand{\mb}{Anamta Akmal}
\newcommand{\mc}{Jagtar Singh}

\newcolumntype{L}[1]{>{\raggedright\arraybackslash}p{#1}}

\lstset{
  basicstyle=\ttfamily\footnotesize,
  breaklines=true,
  frame=single,
  backgroundcolor=\color{gray!8},
  columns=fullflexible,
  keepspaces=true
}

\hypersetup{
    colorlinks=true,
    linkcolor=blue,
    filecolor=magenta,
    urlcolor=cyan,
    citecolor=blue
}

\setlength{\emergencystretch}{3em}

\pagestyle{fancy}
\fancyhf{}
\rhead{SE423 - Milestone 1}
\lhead{Mini API Gateway (MAG)}
\rfoot{Page \thepage}
\setlength{\headheight}{14.5pt}

\titleformat{\section}
 {\normalfont\Large\bfseries\color{black}}{\thesection}{1em}{}
\titleformat{\subsection}
 {\normalfont\large\bfseries\color{black}}{\thesubsection}{1em}{}
\titleformat{\subsubsection}
 {\normalfont\normalsize\bfseries\color{black}}{\thesubsubsection}{1em}{}

\setlength{\parskip}{0.7em}
\setlength{\parindent}{0pt}
\renewcommand{\arraystretch}{1.25}

% ---- UML diagram styles and helpers ----
\newcommand{\st}[1]{\guillemotleft#1\guillemotright}      % stereotype «...»
\newcommand{\cn}[1]{{\centering #1\par}}   % centred class-name compartment

\tikzset{
  % arrows
  assoc/.style={-{Straight Barb[length=5pt,width=6pt]}, line width=0.6pt},      % directed association (open arrow)
  dep/.style={dashed, -{Straight Barb[length=5pt,width=6pt]}, line width=0.6pt},% dependency (dashed, open arrow)
  realize/.style={dashed, -{Triangle[open,length=8pt,width=8pt]}, line width=0.6pt}, % realization
  pl/.style={dashed, line width=0.6pt},                                          % realization feeder line
  call/.style={-{Stealth[length=5pt,width=6pt]}, line width=0.8pt},              % external call
  lbl/.style={font=\tiny, inner sep=1pt, fill=white},
  role/.style={font=\tiny, inner sep=1.5pt, align=left},
  % class box: name | attributes | operations
  class/.style={draw, rectangle split, rectangle split parts=3,
    rectangle split part fill={gray!12,white,white},
    font=\scriptsize, align=left, text width=4.6cm, inner xsep=4pt, inner ysep=3pt},
  note/.style={draw, dashed, font=\tiny, align=left, inner sep=3pt, fill=yellow!10},
  % component diagram
  comp/.style={draw, fill=white, minimum width=3.8cm, minimum height=0.9cm,
    align=center, font=\small, inner xsep=14pt},
  boundary/.style={draw, rounded corners=2pt, fill=gray!5, thick}
}

% component icon (rectangle with two small rectangles on the left edge), top-right corner
\newcommand{\compicon}[1]{%
  \begin{scope}[shift={($(#1.north east)+(-0.32,-0.12)$)}, line width=0.3pt]
    \draw[fill=white] (0,0) rectangle (0.24,-0.16);
    \draw[fill=white] (-0.06,-0.03) rectangle (0.04,-0.07);
    \draw[fill=white] (-0.06,-0.09) rectangle (0.04,-0.13);
  \end{scope}}


\begin{document}

\begin{titlepage}
    \centering
    \includegraphics[width=0.4\textwidth]{Giki Logo.png}
    \vspace{1.0cm}

    {\Huge\bfseries Mini API Gateway (MAG): A Lightweight, Resilient Reverse Proxy \par}
    \vspace{1cm}
    {\LARGE Joint Semester Project -- Milestone 1 \\ Proposal, OOD \& Pattern Plan \par}
    \vspace{1cm}

    {\Large\bfseries SE423: Software Construction and Development \par}
    \vspace{0.5cm}

    {\large Instructor: Dr. Mian Muaz Razaq \par}
    \vspace{0.5cm}

    {\large\bfseries Submitted By: \par}
    \vspace{0.5cm}
    \begin{tabular}{l >{\quad} l}
        \ma & 2023085 \\
        \mb & 2023113 \\
        \mc & 2023266 \\
    \end{tabular}

    \vspace{0.8cm}
    {\large Repository: \code{https://github.com/jagtar5/MiniApiGateway} \par}

    \vspace{\fill}
    \textit{Faculty of Computer Science and Engineering}\\
    \textit{GIK Institute of Engineering Sciences and Technology}\\
    \vspace{0.5cm}
    {\large Date: 11 October 2026 \par}
\end{titlepage}
\thispagestyle{empty}

\newpage
\tableofcontents
\newpage

%====================================================================
\section{Introduction and Rubric Coverage}
%====================================================================
This is the SE423 Milestone~1 submission. It defines the problem, maps each Section~3.1 requirement to a plan, assigns modules to owners, and presents the object-oriented design, SOLID analysis, pattern choices and repository setup.

\begin{table}[H]
\centering
\small
\begin{tabular}{L{6.2cm} c L{5cm}}
\toprule
\textbf{M1 Rubric Criterion} & \textbf{Weight} & \textbf{Where Addressed} \\
\midrule
Problem definition \& requirement mapping & 15 & Sections~\ref{sec:problem}, \ref{sec:mapping} \\
Module breakdown \& ownership & 15 & Section~\ref{sec:modules} \\
Object-oriented design & 25 & Section~\ref{sec:ood} \\
SOLID application plan & 20 & Section~\ref{sec:solid} \\
Design-pattern selection \& justification & 15 & Section~\ref{sec:patterns} \\
Repository \& documentation setup & 10 & Section~\ref{sec:repo} \\
\bottomrule
\end{tabular}
\caption{Rubric coverage.}
\end{table}

%====================================================================
\section{Problem Definition and System Description}
\label{sec:problem}
%====================================================================
\textbf{Mini API Gateway (MAG)} is a lightweight, configuration-driven reverse proxy written in Java with Spring Boot. It sits between clients and upstream HTTP services and applies the following through a filter chain we implement ourselves:

\begin{itemize}[noitemsep]
    \item request identification (\cls{X-Request-Id}) and structured logging;
    \item stateless JWT validation;
    \item concurrent, pluggable rate limiting (token bucket, fixed window, sliding-window counter);
    \item safe caching of GET responses (TTL, LRU eviction, single-flight on miss);
    \item circuit breaking and timeouts for upstream calls;
    \item routing and forwarding;
    \item PostgreSQL-backed, hot-reloadable configuration through an Admin REST API;
    \item a simple demo dashboard (Section~\ref{sec:scope}).
\end{itemize}

MAG is \textbf{generic}: routes, limits and cache rules are data, not code, so it can protect any HTTP service without modification.

\subsection{Users and Stakeholders}
\begingroup
\small
\setlength{\tabcolsep}{4pt}
\begin{tabular}{L{3.4cm} L{5.4cm} L{5.4cm}}
\toprule
\textbf{User} & \textbf{Interaction} & \textbf{Primary need} \\
\midrule
API client & Sends requests to the gateway instead of the backend & Fast, correct responses; clear error codes (401, 429, 503) \\
Platform engineer & Changes routes and limits through the Admin API or dashboard & Runtime changes without redeploying or dropping traffic \\
Backend owner & Places a service behind the gateway & Protection from abuse, overload and cascading failure \\
Course evaluators & Watch behavior in the dashboard; inspect code and tests & Visible, verifiable behavior \\
\bottomrule
\end{tabular}
\captionof{table}{System users and needs.}
\endgroup

\subsection{Scope}
\label{sec:scope}
\textbf{In scope:}
\begin{itemize}[noitemsep]
    \item the five modules of Section~\ref{sec:modules};
    \item test fixtures: dummy upstream services (\cls{/fast}, \cls{/slow}, \cls{/flaky}, \cls{/cacheable}), a JWT-minting utility and a load script;
    \item a \textbf{demo dashboard}: one static HTML/JavaScript page served by the gateway, with three panels: (1) live metrics (allowed/blocked counts, cache hit ratio, circuit state); (2) routes and rules (view, edit a limit through the Admin API); (3) a request tester (send a request with a chosen token, show status and headers).
\end{itemize}

\textbf{Out of scope:}
\begin{itemize}[noitemsep]
    \item reading or modifying request/response bodies (only URLs, headers, methods);
    \item user registration (JWTs are minted elsewhere and only verified here);
    \item distributed mode (documented, not implemented);
    \item dashboard features beyond the three panels: no framework, accounts or styling effort; it contains no logic of its own and only calls the Admin API and metrics endpoint. The admin token is typed into the page and never stored in code.
\end{itemize}

%====================================================================
\section{Requirement-to-Minimum Mapping (Section 3.1)}
\label{sec:mapping}
%====================================================================
\begingroup
\small
\setlength{\tabcolsep}{4pt}
\begin{longtable}{L{3.6cm} L{6.6cm} L{3.6cm}}
\caption{Requirement-to-minimum mapping.}\\
\toprule
\textbf{Requirement} & \textbf{How MAG meets it} & \textbf{Evidence / milestone} \\
\midrule
\endfirsthead
\toprule
\textbf{Requirement} & \textbf{How MAG meets it} & \textbf{Evidence / milestone} \\
\midrule
\endhead
4+ modules with real logic & Five modules: pipeline, identity and traffic control, cache, resilience, configuration and admin (Section~\ref{sec:modules}) & M1: design. M2--M3: code \\
OOD with SOLID justification & All five principles applied with concrete decisions (Section~\ref{sec:solid}) & M1: analysis. M2--M3: verified in code \\
3+ patterns incl.\ creational and structural/behavioral & Six patterns: Chain of Responsibility, Strategy, State (behavioral); Factory, Builder (creational); Proxy (structural) & M1: rationale. M2: running \\
Concurrency with thread-safety strategy & Atomic counters, copy-on-write config, single-flight cache, atomic breaker transitions (Section~\ref{sec:concurrency}) & M1: strategy. M2: stress test \\
Defensive programming, structured exceptions & Exception hierarchy mapped to HTTP codes by one central mapper; boundary validation (Section~\ref{sec:defensive}) & M1: design. M2: code \\
Persistent layer with schema & PostgreSQL with Flyway migrations (Section~\ref{sec:schema}) & M1: schema. M2: loaded at startup \\
TDD, $\geq$60\% coverage, red-green-refactor history & JUnit~5, Mockito, JaCoCo; commit prefixes \cls{test:}, \cls{feat:}, \cls{refactor:}; no squash merges & M1: policy. M3: $\geq$60\% \\
Hosted Git repo, attributable history & GitHub, feature branch per owner, reviewed pull requests & M1: live \\
CI on push & GitHub Actions: build, test, coverage, static analysis & M1: running \\
\bottomrule
\end{longtable}
\endgroup

%====================================================================
\section{Module Breakdown and Ownership}
\label{sec:modules}
%====================================================================
\begingroup
\small
\setlength{\tabcolsep}{4pt}
\begin{longtable}{L{1.0cm} L{3.6cm} L{7.4cm} L{2.2cm}}
\caption{Modules, scope and primary owner.}\\
\toprule
\textbf{ID} & \textbf{Module} & \textbf{One-line scope} & \textbf{Owner} \\
\midrule
\endfirsthead
\toprule
\textbf{ID} & \textbf{Module} & \textbf{One-line scope} & \textbf{Owner} \\
\midrule
\endhead
M-A & Proxy Pipeline & Route matching, our own \cls{GatewayFilter} chain, forwarding, header handling, timeouts, central exception-to-HTTP mapping. & \ma \\
M-B & Identity \& Traffic Control & JWT validation, pluggable client-key extraction, thread-safe rate limiting with three algorithms. & \mb \\
M-C & Fast-Path Cache & Safe GET caching with TTL, bounded LRU eviction and single-flight loading. & \mc \\
M-D & Resilience Engine & Circuit-breaker state machine (closed, open, half-open) with failure tracking and probing. & \mc \\
M-E & Configuration \& Admin & PostgreSQL routes and rules, validation, immutable snapshot hot reload, last-known-good fallback, Admin REST API. & \mb \\
\bottomrule
\end{longtable}
\endgroup

\textbf{Supporting components (not graded modules):} \cls{Clock} abstraction, metrics and logging, dummy upstreams, JWT minter, load script, CI: \ma. Demo dashboard: \mc.

\textbf{Ownership rules:}
\begin{itemize}[noitemsep]
    \item Each member writes, tests and defends their own modules in the demo and Q\&A.
    \item \ma\ delivers the \cls{GatewayFilter} interface and a forward-only walking skeleton early; others build against the interface until then.
    \item Every pull request is reviewed by another member; concurrency code gets a second human review.
\end{itemize}

%====================================================================
\section{Object-Oriented Design}
\label{sec:ood}
%====================================================================

\subsection{Architectural Approach}
MAG is split into a \textbf{plain-Java core} (no Spring annotations) and a thin \textbf{Spring Boot shell}. The core holds all business logic. The shell provides the servlet entry point, dependency-injection wiring, JPA repositories, metrics, and the static dashboard files. The entry point hands each request to \emph{our} filter chain, not Spring's. A multi-module Maven build makes the compiler enforce this: the core cannot import Spring.

\subsection{Component Diagram and Request Flow}
The cache is a proxy behind the \cls{UpstreamClient} interface and the circuit breaker wraps the HTTP client, so a cache hit never reaches the breaker and cannot distort failure statistics.

\begin{figure}[H]
\centering
\resizebox{\linewidth}{!}{%
\begin{tikzpicture}
% ----- actor -----
\begin{scope}[shift={(0,4.3)}]
  \draw (0,0.32) circle (0.15);
  \draw (0,0.17)--(0,-0.25) (-0.22,0.02)--(0.22,0.02) (0,-0.25)--(-0.15,-0.58) (0,-0.25)--(0.15,-0.58);
\end{scope}
\node[anchor=west, font=\small] at (0.35,4.3) {Client / Load tool};

% ----- request pipeline (centre column) -----
\node[comp] (entry) at (0,0)     {Servlet Entry Point\\[-2pt]{\scriptsize (Spring shell)}};
\node[comp] (rid)   at (0,-1.7)  {RequestIdFilter};
\node[comp] (auth)  at (0,-3.4)  {AuthFilter\\[-2pt]{\scriptsize (JWT)}};
\node[comp] (rl)    at (0,-5.1)  {RateLimitFilter};
\node[comp] (fwd)   at (0,-6.8)  {ForwardFilter};
\node[comp] (cache) at (0,-8.5)  {CachingUpstreamClient\\[-2pt]{\scriptsize (cache + single-flight)}};
\node[comp] (res)   at (0,-10.2) {ResilientUpstreamClient\\[-2pt]{\scriptsize (circuit breaker + timeouts)}};
\node[comp] (http)  at (0,-11.9) {HttpUpstreamClient};

% ----- left column -----
\node[comp] (err)   at (-5.6,-3.4) {Central Exception\\Mapper};
\node[comp] (clock) at (-5.6,-9.35) {Clock\\[-2pt]{\scriptsize (injected)}};

% ----- right column -----
\node[comp] (cfg)  at (5.7,-5.1)  {ConfigSnapshot\\[-2pt]{\scriptsize (AtomicReference)}};
\node[comp] (mgr)  at (5.7,-7.4)  {Config Manager\\+ Admin API};
\node[comp] (dash) at (5.7,-9.7)  {Demo Dashboard\\[-2pt]{\scriptsize (static web UI)}};

% ----- external services -----
\node[cylinder, shape border rotate=90, aspect=0.25, draw, fill=white, minimum height=1.6cm,
      minimum width=2.6cm, align=center, font=\small] (pg) at (11.8,-7.4) {PostgreSQL};
\node[comp, minimum width=3.2cm] (up) at (11.8,-11.9) {Upstream\\services};

\foreach \n in {entry,rid,auth,rl,fwd,cache,res,http,err,clock,cfg,mgr,dash,up} {\compicon{\n}}

% ----- boundaries -----
\begin{scope}[on background layer]
  \node[boundary, fit=(entry)(rid)(http)(err)(clock)(cfg)(mgr)(dash),
        inner xsep=0.5cm, inner ysep=0.5cm, yshift=0.4cm] (app) {};
  \node[boundary, fit=(pg)(up), inner xsep=0.5cm, inner ysep=0.5cm, yshift=0.0cm] (ext) {};
\end{scope}
\node[anchor=north west, font=\small\bfseries, inner sep=4pt] at (app.north west) {mini\_api\_gateway (application)};
\node[anchor=north west, font=\small\bfseries, inner sep=4pt] at (ext.north west) {External services};

% ----- actor uses the entry point -----
\draw[assoc] (0,3.65) -- node[right, font=\scriptsize]{uses} (entry.north);

% ----- dependencies inside the application (dashed, <<use>>) -----
\foreach \a/\b in {entry/rid, rid/auth, auth/rl, rl/fwd, fwd/cache, cache/res, res/http}
  \draw[dep] (\a) -- node[right, lbl, fill=none] {\st{use}} (\b);
\draw[dep] (auth) -- node[above, lbl] {\st{use}} (err);
\draw[dep] (rl)   -- node[above, lbl] {\st{use}} (cfg);
\draw[dep] (mgr)  -- node[right, lbl, fill=none] {\st{use}} (cfg);
\draw[dep] (dash) -- node[right, lbl, fill=none] {\st{use} (REST)} (mgr);
\draw[dep] (cache.west) -| node[pos=0.25, above, lbl] {\st{use}} (clock.north);
\draw[dep] (res.west)   -| node[pos=0.25, above, lbl] {\st{use}} (clock.south);

% ----- external calls (solid, labelled with protocol) -----
\draw[call] (mgr.east)  -- node[above, lbl] {SQL}  (pg.west);
\draw[call] (http.east) -- node[above, lbl] {HTTP} (up.west);
\end{tikzpicture}}
\caption{Component and request-flow diagram. Dashed arrows are \st{use} dependencies; solid arrows are calls that cross the application boundary.}
\label{fig:component}
\end{figure}

\newpage
\subsection{Class Diagrams}

\subsubsection{Pipeline, Chain of Responsibility and Builder (M-A)}
\begin{figure}[H]
\centering
\resizebox{\linewidth}{!}{%
\begin{tikzpicture}
\node[class] (gf) at (0,0) {\cn{\st{interface}\\\textbf{GatewayFilter}}\nodepart{second}\strut\nodepart{third}+doFilter(ctx, chain): void};

\node[class] (rid)  at (-8.4,-5) {\cn{\textbf{RequestIdFilter}}\nodepart{second}\strut\nodepart{third}+doFilter(ctx, chain): void};
\node[class] (auth) at (-2.8,-5) {\cn{\textbf{AuthFilter}}\nodepart{second}-tokenVerifier: TokenVerifier\nodepart{third}+doFilter(ctx, chain): void};
\node[class] (rl)   at (2.8,-5)  {\cn{\textbf{RateLimitFilter}}\nodepart{second}-limiterRegistry: LimiterRegistry\\-keyExtractor: KeyExtractor\nodepart{third}+doFilter(ctx, chain): void};
\node[class] (ff)   at (8.4,-5)  {\cn{\textbf{ForwardFilter}}\nodepart{second}-upstreamClient: UpstreamClient\nodepart{third}+doFilter(ctx, chain): void};

\node[class] (fc) at (0,5) {\cn{\textbf{FilterChain}}\nodepart{second}-position: int\nodepart{third}+proceed(ctx): void};
\node[class, text width=5.4cm] (pb) at (-7.6,5) {\cn{\textbf{PipelineBuilder}}\nodepart{second}-routeConfig: RouteConfig\nodepart{third}+withAuth(): PipelineBuilder\\+withRateLimit(): PipelineBuilder\\+build(): FilterChain};
\node[class, text width=5.4cm] (ctx) at (8.4,0) {\cn{\textbf{GatewayContext}}\nodepart{second}-request: GatewayRequest \{readOnly\}\\-response: GatewayResponse\\-attributes: Map<String, Object>\nodepart{third}+shortCircuit(resp): void};

% realization (shared hollow triangle)
\coordinate (jun) at ($(gf.south)+(0,-1.2)$);
\foreach \c in {rid,auth,rl,ff} \draw[pl] (\c.north) |- (jun);
\draw[realize] (jun) -- (gf.south);

% associations / dependencies
\draw[assoc] (fc.south) -- (gf.north)
  node[role, near start, right, xshift=2pt] {-filters\\{*}}
  node[role, near end, right, xshift=2pt, yshift=2pt] {};
\draw[dep] (pb.east) -- node[above, lbl] {\st{create}} (fc.west);
\draw[dep] (gf.east) -- node[above, lbl] {\st{use}} (ctx.west);
\end{tikzpicture}}
\caption{Filter chain and pipeline builder.}
\end{figure}

\subsubsection{Upstream Access, Proxy and Resilience (M-A, M-C, M-D)}
\begin{figure}[H]
\centering
\resizebox{\linewidth}{!}{%
\begin{tikzpicture}
\node[class, text width=6.8cm] (uc) at (0,0) {\cn{\st{interface}\\\textbf{UpstreamClient}}\nodepart{second}\strut\nodepart{third}+execute(req): UpstreamResponse};
\node[class] (ff) at (0,5) {\cn{\textbf{ForwardFilter}}\nodepart{second}\strut\nodepart{third}+doFilter(ctx, chain): void};

\node[class] (cache) at (-8.4,-5.5) {\cn{\textbf{CachingUpstreamClient}}\nodepart{second}-cache: BoundedLruCache\\-inFlight: Map<Key, Future>\nodepart{third}+execute(req): UpstreamResponse};
\node[class] (res)   at (0,-5.5)    {\cn{\textbf{ResilientUpstreamClient}}\nodepart{second}-breaker: CircuitBreaker\nodepart{third}+execute(req): UpstreamResponse};
\node[class] (http)  at (8.4,-5.5)  {\cn{\textbf{HttpUpstreamClient}}\nodepart{second}-connectTimeout: Duration\\-readTimeout: Duration\nodepart{third}+execute(req): UpstreamResponse};

% realization
\coordinate (jun) at ($(uc.south)+(0,-1.2)$);
\foreach \c in {cache,res,http} \draw[pl] (\c.north) |- (jun);
\draw[realize] (jun) -- (uc.south);

% ForwardFilter -> UpstreamClient
\draw[assoc] (ff.south) -- (uc.north)
  node[role, near end, right, xshift=2pt] {-client\\1};

% decorators delegate to the next UpstreamClient
\draw[assoc] (cache.west) -- ++(-0.9,0) |- (uc.west)
  node[role, near end, above] {-delegate\\1};
\draw[assoc] (res.east) -- ++(1.8,0) |- (uc.east)
  node[role, pos=0.62, right] {-delegate\\1};
\end{tikzpicture}}
\caption{\cls{UpstreamClient} proxy chain (cache $\rightarrow$ breaker $\rightarrow$ HTTP).}
\end{figure}

\subsubsection{Rate Limiting: Strategy and Factory (M-B)}
\begin{figure}[H]
\centering
\resizebox{\linewidth}{!}{%
\begin{tikzpicture}
\node[class, text width=4.8cm] (alg) at (0,0) {\cn{\st{interface}\\\textbf{RateLimitAlgorithm}}\nodepart{second}\strut\nodepart{third}+tryAcquire(key): Decision};

\node[class] (tb) at (-7.5,-5.5) {\cn{\textbf{TokenBucket}}\nodepart{second}-capacity: int\\-refillRate: double\\-state: Map<String, Bucket>\\-clock: Clock\nodepart{third}+tryAcquire(key): Decision};
\node[class] (fw) at (0,-5.5) {\cn{\textbf{FixedWindow}}\nodepart{second}-limit: int\\-window: Duration\\-counters: Map<String, Window>\\-clock: Clock\nodepart{third}+tryAcquire(key): Decision};
\node[class] (sw) at (7.5,-5.5) {\cn{\textbf{SlidingWindowCounter}}\nodepart{second}-limit: int\\-window: Duration\\-state: Map<String, Counts>\\-clock: Clock\nodepart{third}+tryAcquire(key): Decision};

\node[class, text width=5.6cm] (fac) at (-9.5,0) {\cn{\textbf{RateLimiterFactory}}\nodepart{second}-registry: Map<String, Creator>\nodepart{third}+create(ruleConfig): RateLimitAlgorithm};
\node[class] (rlf) at (0,5.5) {\cn{\textbf{RateLimitFilter}}\nodepart{second}\strut\nodepart{third}+doFilter(ctx, chain): void};
\node[class] (ke) at (9.5,5.5) {\cn{\st{interface}\\\textbf{KeyExtractor}}\nodepart{second}\strut\nodepart{third}+extract(ctx): String};
\node[class] (jk) at (7,0.3) {\cn{\textbf{JwtSubjectKeyExtractor}}\nodepart{second}\strut\nodepart{third}+extract(ctx): String};
\node[class] (ik) at (12.3,0.3) {\cn{\textbf{ClientIpKeyExtractor}}\nodepart{second}\strut\nodepart{third}+extract(ctx): String};

% realization: algorithms
\coordinate (jun1) at ($(alg.south)+(0,-1.2)$);
\foreach \c in {tb,fw,sw} \draw[pl] (\c.north) |- (jun1);
\draw[realize] (jun1) -- (alg.south);
% realization: key extractors
\coordinate (jun2) at ($(ke.south)+(0,-1.1)$);
\foreach \c in {jk,ik} \draw[pl] (\c.north) |- (jun2);
\draw[realize] (jun2) -- (ke.south);

% associations
\draw[assoc] (rlf.south) -- (alg.north)
  node[role, near end, right, xshift=2pt] {-limiterRegistry\\{*}};
\draw[assoc] (rlf.east) -- (ke.west)
  node[role, near end, above] {-keyExtractor\\1};
% factory creates algorithms
\draw[dep] (fac.east) -- node[above, lbl] {\st{create}} (alg.west);
\end{tikzpicture}}
\caption{Rate-limit algorithms and key extraction (Strategy), \cls{RateLimiterFactory} (Factory).}
\end{figure}

\subsubsection{Circuit Breaker: State Pattern (M-D)}
\begin{figure}[H]
\centering
\resizebox{\linewidth}{!}{%
\begin{tikzpicture}
\node[class] (cb) at (-6.5,0) {\cn{\textbf{CircuitBreaker}}\nodepart{second}-clock: Clock\\-config: BreakerConfig\nodepart{third}+allowRequest(): boolean\\+onSuccess(): void\\+onFailure(): void};
\node[class] (cs) at (2.5,0) {\cn{\st{interface}\\\textbf{CircuitState}}\nodepart{second}\strut\nodepart{third}+allowRequest(ctx): boolean\\+onSuccess(ctx): void\\+onFailure(ctx): void};

\node[class] (cl) at (-5,-5.5) {\cn{\textbf{ClosedState}}\nodepart{second}-failureWindow: FailureWindow\nodepart{third}+allowRequest(ctx): boolean\\+onSuccess(ctx): void\\+onFailure(ctx): void};
\node[class] (op) at (2.5,-5.5) {\cn{\textbf{OpenState}}\nodepart{second}-openedAt: Instant\nodepart{third}+allowRequest(ctx): boolean\\+onSuccess(ctx): void\\+onFailure(ctx): void};
\node[class] (ho) at (10,-5.5) {\cn{\textbf{HalfOpenState}}\nodepart{second}-probesAllowed: int\nodepart{third}+allowRequest(ctx): boolean\\+onSuccess(ctx): void\\+onFailure(ctx): void};

\node[note, anchor=north] (n1) at ($(cl.south)+(0,-0.8)$) {opens when failures\\reach the threshold};
\node[note, anchor=north] (n2) at ($(op.south)+(0,-0.8)$) {half-opens after\\the open timeout};
\node[note, anchor=north] (n3) at ($(ho.south)+(0,-0.8)$) {closes on success,\\re-opens on failure};
\foreach \s/\n in {cl/n1,op/n2,ho/n3} \draw[dashed, line width=0.4pt] (\s.south) -- (\n.north);

% realization
\coordinate (jun) at ($(cs.south)+(0,-1.2)$);
\foreach \c in {cl,op,ho} \draw[pl] (\c.north) |- (jun);
\draw[realize] (jun) -- (cs.south);

% CircuitBreaker delegates to its current state
\draw[assoc] (cb.east) -- (cs.west)
  node[role, near end, above] {-state\\1};
\end{tikzpicture}}
\caption{Circuit breaker (State pattern).}
\end{figure}

\subsection{Abstraction and Encapsulation}
\begin{itemize}[noitemsep]
    \item \textbf{Contracts:} callers depend on small interfaces (\cls{GatewayFilter}, \cls{RateLimitAlgorithm}, \cls{KeyExtractor}, \cls{UpstreamClient}, \cls{CircuitState}, \cls{ConfigProvider}), never on concrete classes.
    \item \textbf{Immutable values:} \cls{GatewayRequest}, \cls{RateLimitDecision}, \cls{ConfigSnapshot} and \cls{RouteConfig} are immutable, which also supports thread safety.
    \item \textbf{Hidden state:} counters, cache entries and breaker counts are private to their owner and reachable only through interface methods; no mutable collection is exposed.
    \item \textbf{Framework isolation:} Spring types stay in the shell; it converts servlet requests to \cls{GatewayRequest} at the boundary.
\end{itemize}

\subsection{Cohesion and Coupling}
\textbf{Cohesion} is functional: each module answers one question. M-B: \emph{may this caller proceed?} M-C: \emph{can this response be reused safely?} M-D: \emph{is this upstream healthy?} M-E: \emph{what are the current rules and how do they change safely?} M-A only orchestrates and holds no policy. Inside M-B, JWT validation, key extraction and quota enforcement are separate classes so security and throughput logic do not mix.

\textbf{Coupling} is kept to interface/data coupling:

\begingroup
\small
\setlength{\tabcolsep}{4pt}
\begin{tabular}{L{2.8cm} L{4.4cm} L{2.2cm} L{4.4cm}}
\toprule
\textbf{Module} & \textbf{Depends on} & \textbf{Coupling} & \textbf{Why acceptable} \\
\midrule
M-A Pipeline & \cls{GatewayFilter}, \cls{UpstreamClient}, \cls{ConfigProvider} & Interface & Knows no limiter, cache or breaker internals \\
M-B Identity \& Traffic & \cls{GatewayContext}, \cls{ConfigProvider}, \cls{Clock} & Interface, data & No dependency on cache, breaker or database \\
M-C Cache & \cls{UpstreamClient}, \cls{Clock} & Interface & Wraps the interface it implements \\
M-D Resilience & \cls{UpstreamClient}, \cls{Clock} & Interface & Independent of cache and limiter \\
M-E Config \& Admin & Repository interface, validator & Interface & Publishes snapshots; calls no other module \\
Dashboard & Admin REST API, metrics endpoint & HTTP only & No access to core classes \\
\bottomrule
\end{tabular}
\captionof{table}{Module dependencies and coupling.}
\endgroup

\textbf{Dependency rules:} (1) the shell depends on the core, never the reverse; (2) modules meet only at interfaces; (3) no cyclic dependencies; (4) time comes only from the injected \cls{Clock}.

%====================================================================
\section{SOLID Application Plan}
\label{sec:solid}
%====================================================================
All five principles are relevant. LSP is only partly observable until M2, when at least two implementations exist.

\begingroup
\small
\setlength{\tabcolsep}{4pt}
\begin{longtable}{L{2.9cm} L{6.6cm} L{4.3cm}}
\caption{SOLID analysis (all five principles demonstrated through concrete decisions).}\\
\toprule
\textbf{Principle / status} & \textbf{Concrete design decision} & \textbf{Justification} \\
\midrule
\endfirsthead
\toprule
\textbf{Principle / status} & \textbf{Concrete design decision} & \textbf{Justification} \\
\midrule
\endhead
\textbf{S} Single Responsibility \newline Relevant & One concern per class: each filter does one job; \cls{ConfigValidator}, \cls{ConfigRepository} and \cls{SnapshotPublisher} are separate; error mapping lives only in the central mapper. & A JWT-rule change must not touch throttling or forwarding code. \\
\textbf{O} Open/ Closed \newline Relevant & A new algorithm is one class plus one factory registration; a new filter is one class plus a builder entry; route behavior comes from configuration. & Extension without editing tested code; verified by an OCP test in M2. \\
\textbf{L} Liskov Substitution \newline Relevant; partly not yet observable (needs $\geq$2 implementations, M2) & One parametrized contract suite runs on \emph{every} \cls{RateLimitAlgorithm} (never exceeds limit, never negative, no exception on valid input). The three \cls{UpstreamClient} implementations must be interchangeable. & Substitutability is provable only with real implementations; the contract is fixed now so the test precedes the second one. \\
\textbf{I} Interface Segregation \newline Relevant & Filters see only read-only \cls{ConfigProvider}; only the admin module sees \cls{ConfigPublisher}. \cls{Clock} exposes only \cls{now()}. Core interfaces have one to three methods. & A filter must not be able to change configuration; small interfaces keep test doubles small. \\
\textbf{D} Dependency Inversion \newline Relevant & The core depends on abstractions (\cls{Clock}, \cls{UpstreamClient}, \cls{ConfigProvider}, \cls{ConfigRepository}); wiring happens only in the shell by constructor injection; the core has no Spring dependency. & Deterministic time-based tests (no sleeps); swappable database and HTTP client. \\
\bottomrule
\end{longtable}
\endgroup

%====================================================================
\section{Design-Pattern Selection and Justification}
\label{sec:patterns}
%====================================================================
Each pattern solves a stated problem and has a proof test; a pattern that fails its proof test is dropped. Six are planned: two creational, one structural, three behavioral.

\begingroup
\small
\setlength{\tabcolsep}{4pt}
\begin{longtable}{L{2.9cm} L{3.5cm} L{3.5cm} L{3.9cm}}
\caption{Patterns with problem, rejected alternative and proof.}\\
\toprule
\textbf{Pattern / location} & \textbf{Problem solved} & \textbf{Rejected alternative} & \textbf{Proof it is genuine} \\
\midrule
\endfirsthead
\toprule
\textbf{Pattern / location} & \textbf{Problem solved} & \textbf{Rejected alternative} & \textbf{Proof it is genuine} \\
\midrule
\endhead
\textbf{Chain of Responsibility} (behavioral) \newline \cls{GatewayFilter} pipeline &
Filters must be orderable, individually testable, and able to stop the request (failed auth skips the rest). &
One large \cls{handle()} method: mixes concerns, cannot be reordered per route. &
A filter is added in a test with no other edits; a short-circuit test shows later filters never run. \\
\textbf{Strategy} (behavioral) \newline \cls{RateLimitAlgorithm}, \cls{KeyExtractor} &
Routes need different algorithms and client identities (JWT subject or IP). &
\cls{switch} on an enum: each new algorithm edits tested code (breaks OCP). &
One contract test runs on every algorithm; a test shows the fixed-window edge burst that token bucket avoids. \\
\textbf{State} (behavioral) \newline \cls{CircuitBreaker} &
Each state reacts differently to allow, success and failure; transition rules would scatter. &
Boolean flags with nested \cls{if}: transitions are hard to read and verify. &
Each state is tested alone; transitions are driven by the injected \cls{Clock} and checked against the transition table. \\
\textbf{Factory} (creational) \newline \cls{RateLimiterFactory} &
Limiters are built from database config at startup and on each reload; callers must not know concrete classes. &
Direct \cls{TokenBucket} construction in filters: ties callers to concrete types and blocks hot reload. &
A new algorithm needs one class and one registry entry, no edits elsewhere (OCP test). \\
\textbf{Builder} (creational) \newline \cls{PipelineBuilder} &
Each route needs a different chain (cache or auth only on some routes). &
One hardcoded chain with per-route \cls{if} checks: unreadable and unsafe to change. &
Two routes from config produce two different chains, verified by comparing filter sequences. \\
\textbf{Proxy} (structural) \newline \cls{CachingUpstreamClient} &
Add caching transparently so no other class knows whether it hit the cache. &
Cache logic inside the forwarder: mixes caching with networking and counts cached hits in breaker statistics. &
Same \cls{UpstreamClient} interface, swappable in tests; a test asserts a cache hit never reaches the breaker. \\
\bottomrule
\end{longtable}
\endgroup

%====================================================================
\section{Supporting Design Decisions}
%====================================================================

\subsection{Concurrency and Thread-Safety Strategy}
\label{sec:concurrency}
\begingroup
\small
\setlength{\tabcolsep}{4pt}
\begin{tabular}{L{2.6cm} L{3.6cm} L{4.4cm} L{3.2cm}}
\toprule
\textbf{Component} & \textbf{Hazard} & \textbf{Strategy} & \textbf{Verification} \\
\midrule
Rate-limit state & Check-then-act race (two threads both see 99/100) & Per-key atomic update with \cls{ConcurrentHashMap.compute} on immutable state & Multi-thread stress test: zero over-limit \\
Configuration & Reader/writer conflict on reload & Immutable \cls{ConfigSnapshot} in an \cls{AtomicReference} (copy-on-write) & Reload under load: no dropped requests \\
Cache & Concurrent eviction; stampede on expiry & Bounded LRU under a lock; per-key single-flight with \cls{CompletableFuture} & 100 concurrent misses give one upstream call \\
Circuit breaker & Concurrent success/failure reports & Atomic counters; state change by compare-and-set & Concurrent failures give one OPEN transition \\
Background work & Idle-key cleanup, config poller & \cls{ScheduledExecutorService} with defined shutdown & Cleanup and shutdown tests \\
\bottomrule
\end{tabular}
\captionof{table}{Thread-safety strategy per component.}
\endgroup

\subsection{Defensive Programming and Exception Handling}
\label{sec:defensive}
\begin{itemize}[noitemsep]
    \item \textbf{Hierarchy:} \cls{GatewayException} base with \cls{AuthenticationException} (401), \cls{AuthorizationException} (403), \cls{RateLimitExceededException} (429), \cls{UpstreamUnavailableException} (502), \cls{CircuitOpenException} (503), \cls{UpstreamTimeoutException} (504) and \cls{InvalidConfigurationException}.
    \item \textbf{One mapping point:} a central mapper converts exceptions to HTTP responses; no generic \cls{catch (Exception)}.
    \item \textbf{Boundary validation:} every module validates its inputs (null, empty, negative limits, malformed JWT, bad route patterns); configuration is validated before it becomes a snapshot.
    \item \textbf{Fail policy:} fail-open versus fail-closed on dependency failure is recorded in an Architecture Decision Record.
\end{itemize}

\subsection{Persistent Data Layer}
\label{sec:schema}
PostgreSQL with versioned Flyway migrations:

\begingroup
\small
\setlength{\tabcolsep}{4pt}
\begin{tabular}{L{3.0cm} L{7.2cm} L{3.6cm}}
\toprule
\textbf{Table} & \textbf{Key columns} & \textbf{Purpose} \\
\midrule
\cls{routes} & \cls{id}, \cls{path_pattern}, \cls{upstream_urls}, \cls{methods}, \cls{auth_required}, \cls{cache_enabled}, \cls{cache_ttl_seconds}, \cls{priority} & Routing rules by priority \\
\cls{rate_limit_rules} & \cls{id}, \cls{route_id} (FK), \cls{algorithm}, \cls{key_strategy}, \cls{limit}, \cls{window_seconds}, \cls{burst} & Per-route limiter settings \\
\cls{config_versions} & \cls{version}, \cls{created_at}, \cls{created_by}, \cls{snapshot_json}, \cls{is_active} & Versioning, last-known-good \\
\bottomrule
\end{tabular}
\captionof{table}{Preliminary schema.}
\endgroup

\subsection{Testing Approach}
\begin{itemize}[noitemsep]
    \item Plain JUnit~5 and Mockito for the core; Spring context tests only for a few integration scenarios.
    \item Injected fake \cls{Clock}, so time-based tests are deterministic with no sleeps.
    \item JaCoCo runs in CI from M1; $\geq$60\% coverage is enforced by M3.
\end{itemize}

%====================================================================
\section{Repository and Documentation Setup}
\label{sec:repo}
%====================================================================
\textbf{Repository:} \code{https://github.com/jagtar5/MiniApiGateway}.

\begin{lstlisting}
mini-api-gateway/
|-- gateway-core/        # plain Java: filters, algorithms, cache, breaker
|-- gateway-app/         # Spring Boot shell: wiring, JPA, admin API,
|                        #   static demo dashboard (resources/static)
|-- dummy-upstreams/     # test fixtures: /fast /slow /flaky /cacheable
|-- tools/               # JWT minter, load scripts
|-- docs/                # design docs, ADRs, AI usage log
|-- .github/workflows/   # CI pipeline
|-- README.md
\end{lstlisting}

\subsection{Branching Strategy}
A lightweight \emph{GitHub Flow} variant for a three-person team:

\begingroup
\small
\setlength{\tabcolsep}{4pt}
\begin{tabular}{L{3.0cm} L{10.8cm}}
\toprule
\textbf{Element} & \textbf{Rule} \\
\midrule
\cls{main} & Protected and always releasable; changes only by pull request with green CI. \\
Feature branches & \cls{feature/<owner>-<topic>}, short-lived, one concern each; also \cls{fix/...} and \cls{docs/...}. \\
Pull requests & Required for every merge; reviewed by a non-author; concurrency changes get a second human review. \\
Merge method & Merge commit or rebase-and-merge. \textbf{Squash is disabled} to preserve the red-green-refactor trail. \\
Commit convention & \cls{test:} (red), \cls{feat:} (green), \cls{refactor:}, \cls{fix:}, \cls{docs:}, \cls{chore:}; small commits that build. \\
Baselines & Annotated tags \cls{m1-baseline}, \cls{m2-baseline}, \cls{m3-baseline}. \\
Quality gate & CI (build, tests, coverage, static analysis) must pass before merge. \\
\bottomrule
\end{tabular}
\captionof{table}{Branching, merge and baseline rules.}
\endgroup

\subsection{CI, README and AI Usage Log}
\textbf{CI:} a GitHub Actions workflow runs on every push and pull request: build with Maven, run tests, publish the JaCoCo report, run SpotBugs and Checkstyle. PostgreSQL comes from a service container or Testcontainers.

\textbf{README:} project summary, architecture and request-flow diagram, prerequisites (JDK, Maven, Docker), build/test/run steps, a route and rule configuration example, branching and commit conventions, and the ownership table. \cls{docs/} holds this document, ADRs (fail-open vs.\ fail-closed, sliding-window counter, snapshot swap vs.\ locking) and the AI Usage Log.


\end{document}